%---------------------------------------------------
% First page with book license
% to add it:   \licensepage
%---------------------------------------------------
\newcommand{\licensepage}[0]{

\null
\vfill
\begin{flushleft}
    \includegraphics[]{img/seal.png}

    Copyright \textcopyright \ \@theauthor \hspace{1pt} (\href{mailto:me@somewhere.com}{rugoli@gmail.com})
    \begin{itemize}
        \item Github: \href{https://github.com/yuki}{https://github.com/yuki}
        \item Twitter: \href{https://twitter.com/rugoli}{https://twitter.com/rugoli}
    \end{itemize}


    \vspace{5mm}
    \textbf{Licencia: } \href{https://creativecommons.org/licenses/by-sa/4.0/}{Creative Commons BY-SA 4.0}


    Este libro se ha relizado teniendo en cuenta la cultura libre. Puedes utilizarlo, modificarlo y compartirlo teniendo en cuenta la licencia \href{https://creativecommons.org/licenses/by-sa/4.0/}{Attribution-ShareAlike} de \textbf{Creative Commons}. Es por eso que:

    \begin{itemize}
        \item \textbf{Atribución}: Debes darme crédito de manera adecuada e incluir un enlace a la licencia e indicar si se han realizado cambios.
        \item \textbf{CompartirIgual}: Si reutilizas, modificas o creas a partir de este material, debes distribuir el trabajo bajo la misma licencia.
    \end{itemize}


    Puedes encontrar la última versión de este libro en la siguiente URL (\href{https://github.com/yuki/my-books}{https://github.com/yuki/my-books}), así como otros libros que he creado. También podrás descargar el código fuente en formato \LaTeX{}.

    \infobox{\textbf{Por favor, ponte en contacto conmigo si encuentras algún fallo, falta de ortografía o quieres mejorar de alguna manera este libro. Gracias.}}
    \vspace{5mm}
\end{flushleft}

\clearpage

} % end of /licensepage command